% ================================================================
% 问题三新增符号表
% 使用方法：在 main.tex 的 longtable 符号表中，删除原来的
% “问题三决策变量”小节（x_r、y_{ru}、z_p、T_0 四行），
% 然后在同一位置写入：% ================================================================
% 问题三新增符号表
% 使用方法：在 main.tex 的 longtable 符号表中，删除原来的
% “问题三决策变量”小节（x_r、y_{ru}、z_p、T_0 四行），
% 然后在同一位置写入：% ================================================================
% 问题三新增符号表
% 使用方法：在 main.tex 的 longtable 符号表中，删除原来的
% “问题三决策变量”小节（x_r、y_{ru}、z_p、T_0 四行），
% 然后在同一位置写入：% ================================================================
% 问题三新增符号表
% 使用方法：在 main.tex 的 longtable 符号表中，删除原来的
% “问题三决策变量”小节（x_r、y_{ru}、z_p、T_0 四行），
% 然后在同一位置写入：\input{q3_symbols_addition.tex}
% 本文件只包含 longtable 内部行，不能作为独立文档编译。
% ================================================================

\midrule
\multicolumn{2}{l}{\textbf{问题三新增集合、参数与决策变量}} \\[2pt]

$\mathcal{D}$
& 问题三的规划日集合，$\mathcal{D}=\{0,1,\ldots,6\}$，分别对应
2026年8月3日至8月9日。 \\

$q$
& 规划日索引，$q\in\mathcal{D}$。 \\

$\Omega_q$
& 第 $q$ 日相对规划起点的分钟偏移，$\Omega_q=1440q$。 \\

$P_{\mathrm{req}}$
& 必须完成运输的非临时人员集合，
$P_{\mathrm{req}}=P\setminus P_{\mathrm{temp}}$。 \\

$\mathcal{C}$
& 路线层生成的时间可行候选架次集合。 \\

$P_r$
& 候选架次 $r$ 计划运输的人员集合。 \\

$\mathcal{C}_p$
& 能够运输人员 $p$ 的候选架次集合，
$\mathcal{C}_p=\{r\in\mathcal{C}:p\in P_r\}$。 \\

$q(r)$
& 候选架次 $r$ 所属的规划日。 \\

$a(u),k(u)$
& 具体飞机 $u$ 的所属机场和机型。 \\

$H_r$
& 能够执行候选架次 $r$ 的具体飞机集合，
$H_r=\{u\in H:a(u)=a(r),\,k(u)=k(r)\}$。 \\

$x_r$
& 候选架次选择变量；若选择候选架次 $r$，则 $x_r=1$，
否则 $x_r=0$。 \\

$y_{ru}$
& 架次--飞机分配变量；若候选架次 $r$ 由具体飞机 $u$ 执行，
则 $y_{ru}=1$，否则 $y_{ru}=0$。 \\

$T_{rh}^{\mathrm{arr}}$
& 架次 $r$ 到达第 $h$ 个停靠点的绝对分钟时刻，以
2026年8月3日00:00为时间原点。 \\

$T_{rh}^{\mathrm{dep}}$
& 架次 $r$ 离开第 $h$ 个停靠点的绝对分钟时刻。 \\

$\pi_{pr}$
& 人员 $p$ 在候选架次 $r$ 上的上机位置序号。 \\

$\delta_{pr}$
& 人员 $p$ 在候选架次 $r$ 上的下机位置序号，满足
$\pi_{pr}<\delta_{pr}$。 \\

$T_{pr}^{\mathrm{pass}}$
& 人员 $p$ 乘坐候选架次 $r$ 时的实际在途时间，
$T_{pr}^{\mathrm{pass}}=
T_{r,\delta_{pr}}^{\mathrm{arr}}-T_{r,\pi_{pr}}^{\mathrm{dep}}$。 \\

$T_r^{\mathrm{air}}$
& 候选架次 $r$ 的实际飞机使用时间，
$T_r^{\mathrm{air}}=T_{r,m_r+1}^{\mathrm{arr}}-T_{r0}^{\mathrm{dep}}$。 \\

$I_{ru}$
& 候选架次 $r$ 分配给飞机 $u$ 时建立的可选时间区间；区间末端
包含30 min机场周转时间，其存在变量为 $y_{ru}$。 \\

$z_p$
& 临时人员服务变量；若临时人员 $p$ 被纳入最终运输方案，
则 $z_p=1$，否则 $z_p=0$。 \\

$S_0$
& 第一阶段完成全部非临时任务后得到的最佳可行基准方案。 \\

$T_0$
& 基准方案 $S_0$ 的总飞机使用时间，
$T_0=T_{\mathrm{air}}(S_0)$，单位为 min。 \\

$M_p$
& 已服务人员 $p$ 的送达时间裕度，
$M_p=l_p-T_{r,\delta_{pr}}^{\mathrm{arr}}$，单位为 min。 \\

$R_{\mathrm{temp}}$
& 临时任务满足率，
$R_{\mathrm{temp}}=
\left(\sum_{p\in P_{\mathrm{temp}}}z_p/160\right)\times100\%$。 \\

% 本文件只包含 longtable 内部行，不能作为独立文档编译。
% ================================================================

\midrule
\multicolumn{2}{l}{\textbf{问题三新增集合、参数与决策变量}} \\[2pt]

$\mathcal{D}$
& 问题三的规划日集合，$\mathcal{D}=\{0,1,\ldots,6\}$，分别对应
2026年8月3日至8月9日。 \\

$q$
& 规划日索引，$q\in\mathcal{D}$。 \\

$\Omega_q$
& 第 $q$ 日相对规划起点的分钟偏移，$\Omega_q=1440q$。 \\

$P_{\mathrm{req}}$
& 必须完成运输的非临时人员集合，
$P_{\mathrm{req}}=P\setminus P_{\mathrm{temp}}$。 \\

$\mathcal{C}$
& 路线层生成的时间可行候选架次集合。 \\

$P_r$
& 候选架次 $r$ 计划运输的人员集合。 \\

$\mathcal{C}_p$
& 能够运输人员 $p$ 的候选架次集合，
$\mathcal{C}_p=\{r\in\mathcal{C}:p\in P_r\}$。 \\

$q(r)$
& 候选架次 $r$ 所属的规划日。 \\

$a(u),k(u)$
& 具体飞机 $u$ 的所属机场和机型。 \\

$H_r$
& 能够执行候选架次 $r$ 的具体飞机集合，
$H_r=\{u\in H:a(u)=a(r),\,k(u)=k(r)\}$。 \\

$x_r$
& 候选架次选择变量；若选择候选架次 $r$，则 $x_r=1$，
否则 $x_r=0$。 \\

$y_{ru}$
& 架次--飞机分配变量；若候选架次 $r$ 由具体飞机 $u$ 执行，
则 $y_{ru}=1$，否则 $y_{ru}=0$。 \\

$T_{rh}^{\mathrm{arr}}$
& 架次 $r$ 到达第 $h$ 个停靠点的绝对分钟时刻，以
2026年8月3日00:00为时间原点。 \\

$T_{rh}^{\mathrm{dep}}$
& 架次 $r$ 离开第 $h$ 个停靠点的绝对分钟时刻。 \\

$\pi_{pr}$
& 人员 $p$ 在候选架次 $r$ 上的上机位置序号。 \\

$\delta_{pr}$
& 人员 $p$ 在候选架次 $r$ 上的下机位置序号，满足
$\pi_{pr}<\delta_{pr}$。 \\

$T_{pr}^{\mathrm{pass}}$
& 人员 $p$ 乘坐候选架次 $r$ 时的实际在途时间，
$T_{pr}^{\mathrm{pass}}=
T_{r,\delta_{pr}}^{\mathrm{arr}}-T_{r,\pi_{pr}}^{\mathrm{dep}}$。 \\

$T_r^{\mathrm{air}}$
& 候选架次 $r$ 的实际飞机使用时间，
$T_r^{\mathrm{air}}=T_{r,m_r+1}^{\mathrm{arr}}-T_{r0}^{\mathrm{dep}}$。 \\

$I_{ru}$
& 候选架次 $r$ 分配给飞机 $u$ 时建立的可选时间区间；区间末端
包含30 min机场周转时间，其存在变量为 $y_{ru}$。 \\

$z_p$
& 临时人员服务变量；若临时人员 $p$ 被纳入最终运输方案，
则 $z_p=1$，否则 $z_p=0$。 \\

$S_0$
& 第一阶段完成全部非临时任务后得到的最佳可行基准方案。 \\

$T_0$
& 基准方案 $S_0$ 的总飞机使用时间，
$T_0=T_{\mathrm{air}}(S_0)$，单位为 min。 \\

$M_p$
& 已服务人员 $p$ 的送达时间裕度，
$M_p=l_p-T_{r,\delta_{pr}}^{\mathrm{arr}}$，单位为 min。 \\

$R_{\mathrm{temp}}$
& 临时任务满足率，
$R_{\mathrm{temp}}=
\left(\sum_{p\in P_{\mathrm{temp}}}z_p/160\right)\times100\%$。 \\

% 本文件只包含 longtable 内部行，不能作为独立文档编译。
% ================================================================

\midrule
\multicolumn{2}{l}{\textbf{问题三新增集合、参数与决策变量}} \\[2pt]

$\mathcal{D}$
& 问题三的规划日集合，$\mathcal{D}=\{0,1,\ldots,6\}$，分别对应
2026年8月3日至8月9日。 \\

$q$
& 规划日索引，$q\in\mathcal{D}$。 \\

$\Omega_q$
& 第 $q$ 日相对规划起点的分钟偏移，$\Omega_q=1440q$。 \\

$P_{\mathrm{req}}$
& 必须完成运输的非临时人员集合，
$P_{\mathrm{req}}=P\setminus P_{\mathrm{temp}}$。 \\

$\mathcal{C}$
& 路线层生成的时间可行候选架次集合。 \\

$P_r$
& 候选架次 $r$ 计划运输的人员集合。 \\

$\mathcal{C}_p$
& 能够运输人员 $p$ 的候选架次集合，
$\mathcal{C}_p=\{r\in\mathcal{C}:p\in P_r\}$。 \\

$q(r)$
& 候选架次 $r$ 所属的规划日。 \\

$a(u),k(u)$
& 具体飞机 $u$ 的所属机场和机型。 \\

$H_r$
& 能够执行候选架次 $r$ 的具体飞机集合，
$H_r=\{u\in H:a(u)=a(r),\,k(u)=k(r)\}$。 \\

$x_r$
& 候选架次选择变量；若选择候选架次 $r$，则 $x_r=1$，
否则 $x_r=0$。 \\

$y_{ru}$
& 架次--飞机分配变量；若候选架次 $r$ 由具体飞机 $u$ 执行，
则 $y_{ru}=1$，否则 $y_{ru}=0$。 \\

$T_{rh}^{\mathrm{arr}}$
& 架次 $r$ 到达第 $h$ 个停靠点的绝对分钟时刻，以
2026年8月3日00:00为时间原点。 \\

$T_{rh}^{\mathrm{dep}}$
& 架次 $r$ 离开第 $h$ 个停靠点的绝对分钟时刻。 \\

$\pi_{pr}$
& 人员 $p$ 在候选架次 $r$ 上的上机位置序号。 \\

$\delta_{pr}$
& 人员 $p$ 在候选架次 $r$ 上的下机位置序号，满足
$\pi_{pr}<\delta_{pr}$。 \\

$T_{pr}^{\mathrm{pass}}$
& 人员 $p$ 乘坐候选架次 $r$ 时的实际在途时间，
$T_{pr}^{\mathrm{pass}}=
T_{r,\delta_{pr}}^{\mathrm{arr}}-T_{r,\pi_{pr}}^{\mathrm{dep}}$。 \\

$T_r^{\mathrm{air}}$
& 候选架次 $r$ 的实际飞机使用时间，
$T_r^{\mathrm{air}}=T_{r,m_r+1}^{\mathrm{arr}}-T_{r0}^{\mathrm{dep}}$。 \\

$I_{ru}$
& 候选架次 $r$ 分配给飞机 $u$ 时建立的可选时间区间；区间末端
包含30 min机场周转时间，其存在变量为 $y_{ru}$。 \\

$z_p$
& 临时人员服务变量；若临时人员 $p$ 被纳入最终运输方案，
则 $z_p=1$，否则 $z_p=0$。 \\

$S_0$
& 第一阶段完成全部非临时任务后得到的最佳可行基准方案。 \\

$T_0$
& 基准方案 $S_0$ 的总飞机使用时间，
$T_0=T_{\mathrm{air}}(S_0)$，单位为 min。 \\

$M_p$
& 已服务人员 $p$ 的送达时间裕度，
$M_p=l_p-T_{r,\delta_{pr}}^{\mathrm{arr}}$，单位为 min。 \\

$R_{\mathrm{temp}}$
& 临时任务满足率，
$R_{\mathrm{temp}}=
\left(\sum_{p\in P_{\mathrm{temp}}}z_p/160\right)\times100\%$。 \\

% 本文件只包含 longtable 内部行，不能作为独立文档编译。
% ================================================================

\midrule
\multicolumn{2}{l}{\textbf{问题三新增集合、参数与决策变量}} \\[2pt]

$\mathcal{D}$
& 问题三的规划日集合，$\mathcal{D}=\{0,1,\ldots,6\}$，分别对应
2026年8月3日至8月9日。 \\

$q$
& 规划日索引，$q\in\mathcal{D}$。 \\

$\Omega_q$
& 第 $q$ 日相对规划起点的分钟偏移，$\Omega_q=1440q$。 \\

$P_{\mathrm{req}}$
& 必须完成运输的非临时人员集合，
$P_{\mathrm{req}}=P\setminus P_{\mathrm{temp}}$。 \\

$\mathcal{C}$
& 路线层生成的时间可行候选架次集合。 \\

$P_r$
& 候选架次 $r$ 计划运输的人员集合。 \\

$\mathcal{C}_p$
& 能够运输人员 $p$ 的候选架次集合，
$\mathcal{C}_p=\{r\in\mathcal{C}:p\in P_r\}$。 \\

$q(r)$
& 候选架次 $r$ 所属的规划日。 \\

$a(u),k(u)$
& 具体飞机 $u$ 的所属机场和机型。 \\

$H_r$
& 能够执行候选架次 $r$ 的具体飞机集合，
$H_r=\{u\in H:a(u)=a(r),\,k(u)=k(r)\}$。 \\

$x_r$
& 候选架次选择变量；若选择候选架次 $r$，则 $x_r=1$，
否则 $x_r=0$。 \\

$y_{ru}$
& 架次--飞机分配变量；若候选架次 $r$ 由具体飞机 $u$ 执行，
则 $y_{ru}=1$，否则 $y_{ru}=0$。 \\

$T_{rh}^{\mathrm{arr}}$
& 架次 $r$ 到达第 $h$ 个停靠点的绝对分钟时刻，以
2026年8月3日00:00为时间原点。 \\

$T_{rh}^{\mathrm{dep}}$
& 架次 $r$ 离开第 $h$ 个停靠点的绝对分钟时刻。 \\

$\pi_{pr}$
& 人员 $p$ 在候选架次 $r$ 上的上机位置序号。 \\

$\delta_{pr}$
& 人员 $p$ 在候选架次 $r$ 上的下机位置序号，满足
$\pi_{pr}<\delta_{pr}$。 \\

$T_{pr}^{\mathrm{pass}}$
& 人员 $p$ 乘坐候选架次 $r$ 时的实际在途时间，
$T_{pr}^{\mathrm{pass}}=
T_{r,\delta_{pr}}^{\mathrm{arr}}-T_{r,\pi_{pr}}^{\mathrm{dep}}$。 \\

$T_r^{\mathrm{air}}$
& 候选架次 $r$ 的实际飞机使用时间，
$T_r^{\mathrm{air}}=T_{r,m_r+1}^{\mathrm{arr}}-T_{r0}^{\mathrm{dep}}$。 \\

$I_{ru}$
& 候选架次 $r$ 分配给飞机 $u$ 时建立的可选时间区间；区间末端
包含30 min机场周转时间，其存在变量为 $y_{ru}$。 \\

$z_p$
& 临时人员服务变量；若临时人员 $p$ 被纳入最终运输方案，
则 $z_p=1$，否则 $z_p=0$。 \\

$S_0$
& 第一阶段完成全部非临时任务后得到的最佳可行基准方案。 \\

$T_0$
& 基准方案 $S_0$ 的总飞机使用时间，
$T_0=T_{\mathrm{air}}(S_0)$，单位为 min。 \\

$M_p$
& 已服务人员 $p$ 的送达时间裕度，
$M_p=l_p-T_{r,\delta_{pr}}^{\mathrm{arr}}$，单位为 min。 \\

$R_{\mathrm{temp}}$
& 临时任务满足率，
$R_{\mathrm{temp}}=
\left(\sum_{p\in P_{\mathrm{temp}}}z_p/160\right)\times100\%$。 \\
